% Gemeinsame Pakete für Vortrag und Ausarbeitung
% Common packages for presentation and paper

%%%%%%%%%%%%%%%%%%%%%%%%%%%%%%%%%%%%%%%%%%%%%%%%%%%

% HINWEIS: natbib wird nicht verwendet, da wir biblatex nutzen (siehe Hauptdokument)
% natbib ist inkompatibel mit biblatex

%%%%%%%%%%%%%%%%%%%%%%%%%%%%%%%%%%%%%%%%%%%%%%%%%%%

% Bildoptionen?
\usepackage{pict2e}

% Latex - Symbole
\usepackage{latexsym}

% Zusätzliche Kodierungen
\usepackage{inconsolata}

% Nicht - umbrechbare Leerzeichen
\usepackage{newunicodechar}

% Abbildungen
\usepackage{tikz}
\usepackage{pgfplots}
\pgfplotsset{compat=1.18}

% Weitere TikZ - Erweiterungen
\usetikzlibrary{chains,shapes.multipart}

% Erweiterungen für die Abbildungen mit Tikz
\usetikzlibrary{shapes,arrows,positioning}

% Abbildungen nebeneinander stellen und Text anordnen
\usepackage{subcaption}

% Mathematische Symbole (dsfont, amssymb)
\usepackage{dsfont}
\usepackage{amssymb}

% Schöner Aufzählungspunkt; und andere schöne Textelemente
\usepackage{textcomp}

% Pfeiloptimierung; und ähnliches
\usepackage{stmaryrd}

% Blindtext (zur Hilfe beim Schreiben)
\usepackage{blindtext}

% Text-Boxen (bei den algorithmischen Problemen)
\usepackage{tcolorbox}

% Grafiken
\usepackage{graphicx}

% Für korrekte Anführungszeichen mit biblatex (behebt Warning)
\usepackage{csquotes}

% Farbiger Text
\usepackage{xcolor}
\usepackage{colortbl}

% Erweiterungen für den Mathemodus
\usepackage{amsmath}

% Erweiterung: Mathefonts
\usepackage{amsfonts}

% Noch mehr Mathe:
\usepackage[centerdots]{mathtools}

% Für Code/Pseudocode
\usepackage{listings}

% ¿Eventuell erforderlich?
\usepackage{array}

% Klickbare Verweise
\usepackage{hyperref}

% Schönere Schriftart
\usepackage{lmodern}

% Tabellen ausrichten
\usepackage{float}

% Mehr (mathematische) Symbole
\usepackage{tipa}

% Um schönere Brüche darstellen zu können
\usepackage{nicefrac}

% um Todo Anmerkungen einfügen zu können
\usepackage{todonotes}
\usepackage{enumitem}
\usepackage{booktabs}

% damit wir Algorithmen in Pseudocode einfach setzen können
\usepackage{algorithm}
\usepackage{algpseudocode}
